% =========================================================================
%  Beat by Bit — Documentação de Incrementos (Semestre 2)
%  Comparação entre labdig_files e labdig_files_midiAudio
% =========================================================================
\documentclass[12pt, a4paper]{article}

% --- Pacotes ---------------------------------------------------------------
\usepackage{fontspec}
\usepackage[brazil]{babel}
\usepackage{geometry}
\usepackage{graphicx}
\usepackage{xcolor}
\usepackage{listings}
\usepackage{booktabs}
\usepackage{tabularx}
\usepackage{longtable}
\usepackage{array}
\usepackage{multirow}
\usepackage{hyperref}
\usepackage{tikz}
\usepackage{amsmath}
\usepackage{fancyhdr}
\usepackage{titlesec}
\usepackage{enumitem}
\usepackage{mdframed}
\usepackage{caption}
\usepackage{subcaption}
\usepackage{amssymb}
% lmodern substituido por fontspec

\usetikzlibrary{automata, positioning, arrows.meta, shapes.geometric, fit, backgrounds}

\geometry{
    top=2.5cm, bottom=2.5cm,
    left=2.5cm, right=2.5cm,
    headheight=15pt
}

% --- Cores -----------------------------------------------------------------
\definecolor{codebg}{RGB}{248,248,248}
\definecolor{codeframe}{RGB}{220,220,220}
\definecolor{vgreen}{RGB}{40,140,40}
\definecolor{vblue}{RGB}{30,80,180}
\definecolor{vred}{RGB}{180,40,40}
\definecolor{vorange}{RGB}{200,100,0}
\definecolor{highlight}{RGB}{255,248,220}
\definecolor{highlightborder}{RGB}{220,180,0}
\definecolor{newcolor}{RGB}{0,120,0}
\definecolor{oldcolor}{RGB}{180,0,0}
\definecolor{sectioncolor}{RGB}{20,60,120}

% --- Estilo de código Verilog ----------------------------------------------
\lstdefinelanguage{Verilog}{
    keywords={module, endmodule, input, output, wire, reg, always, assign,
              begin, end, if, else, case, endcase, posedge, negedge,
              localparam, parameter, initial, default, or},
    keywordstyle=\color{vblue}\bfseries,
    comment=[l]{//},
    commentstyle=\color{vgreen}\itshape,
    morestring=[b]",
    stringstyle=\color{vred},
    sensitive=true
}

\lstset{
    language=Verilog,
    basicstyle=\ttfamily\footnotesize,
    backgroundcolor=\color{codebg},
    frame=single,
    rulecolor=\color{codeframe},
    breaklines=true,
    tabsize=4,
    showstringspaces=false,
    numbers=left,
    numberstyle=\tiny\color{gray},
    numbersep=6pt,
    xleftmargin=1.5em,
    framexleftmargin=1.5em,
    captionpos=b
}

% --- Cabeçalho/Rodapé ------------------------------------------------------
\pagestyle{fancy}
\fancyhf{}
\fancyhead[L]{\small\textit{Beat by Bit — Documentação de Incrementos}}
\fancyhead[R]{\small\textit{Laboratório de Digitais}}
\fancyfoot[C]{\thepage}
\renewcommand{\headrulewidth}{0.4pt}

% --- Formatação de seções --------------------------------------------------
\titleformat{\section}{\large\bfseries\color{sectioncolor}}{\thesection.}{0.6em}{}[\titlerule]
\titleformat{\subsection}{\normalsize\bfseries\color{sectioncolor}}{\thesubsection.}{0.5em}{}
\titleformat{\subsubsection}{\normalsize\bfseries}{\thesubsubsection.}{0.5em}{}

% --- Ambientes personalizados ----------------------------------------------
\newmdenv[
    backgroundcolor=highlight,
    linecolor=highlightborder,
    linewidth=1pt,
    innerleftmargin=8pt,
    innerrightmargin=8pt,
    innertopmargin=6pt,
    innerbottommargin=6pt,
    skipabove=6pt,
    skipbelow=6pt
]{destaque}

\newcommand{\novo}[1]{\textcolor{newcolor}{\textbf{[NOVO]}} #1}
\newcommand{\alterado}[1]{\textcolor{vorange}{\textbf{[ALT]}} #1}
\newcommand{\antes}[1]{\textcolor{oldcolor}{\texttt{#1}}}
\newcommand{\depois}[1]{\textcolor{newcolor}{\texttt{#1}}}
\newcommand{\vfile}[1]{\texttt{\textbf{#1}}}
\newcommand{\vsig}[1]{\texttt{#1}}

% ==========================================================================
\begin{document}

% --- Página de título ------------------------------------------------------
\begin{titlepage}
    \centering
    \vspace*{2cm}

    {\Huge\bfseries\color{sectioncolor} Beat by Bit}\\[0.5em]
    {\large Jogo de Ritmo em FPGA — Placa DE0-CV}\\[2em]

    \rule{\linewidth}{1pt}\\[0.8em]
    {\LARGE\bfseries Documentação de Incrementos}\\[0.4em]
    {\large Semestre 2 — Revisão de Código}\\[0.8em]
    \rule{\linewidth}{1pt}

    \vspace{2em}

    \begin{tabular}{ll}
        \textbf{Versão base:}      & \texttt{labdig\_files/projeto\_sem2} \\[0.3em]
        \textbf{Versão atual:}     & \texttt{labdig\_files\_midiAudio/projeto\_sem2} \\[0.3em]
        \textbf{Clock do sistema:} & 50\,MHz (Cyclone\,V) \\[0.3em]
        \textbf{Taxa de quadros:}  & 60\,fps \\[0.3em]
        \textbf{Baud UART:}        & 115\,200\,bps \\
    \end{tabular}

    \vspace{3em}

    \begin{destaque}
    \centering
    \textbf{Resumo das mudanças principais}\\[0.5em]
    \begin{tabular}{ll}
        ROM do chart: & 4\,428 $\to$ 9\,861 entradas ($\approx$73\,s $\to$ $\approx$164\,s) \\
        Frame counter: & 13\,bits $\to$ 14\,bits \\
        Largura da ROM: & 4\,bits $\to$ 8\,bits por entrada \\
        Pacote UART: & 20 bytes $\to$ 22 bytes \\
        FSM PAUSE: & nova transição \texttt{PAUSE $\to$ IDLE} via \texttt{song\_ok} \\
        Debounce btn0: & reset desacoplado de \texttt{limpaR} \\
    \end{tabular}
    \end{destaque}

    \vfill
    {\small Laboratório de Dispositivos Digitais --- 2025/2026}
\end{titlepage}

% --- Sumário ---------------------------------------------------------------
\tableofcontents
\newpage

% ==========================================================================
\section{Visão Geral do Projeto}
% ==========================================================================

O \textbf{Beat by Bit} é um jogo de ritmo implementado inteiramente em hardware
descritivo Verilog, sintetizado para a FPGA Cyclone\,V da placa Altera DE0-CV.
A lógica do jogo — geração de quadros, leitura do chart, detecção de acertos e
erros, contagem de placar — reside nos módulos RTL. Uma aplicação Python em
PC recebe o estado do jogo quadro a quadro via UART e realiza toda a
renderização gráfica e reprodução de áudio.

\subsection{Estrutura modular}

A hierarquia de módulos é apresentada na Tabela~\ref{tab:hierarquia}.

\begin{table}[h!]
\centering
\caption{Hierarquia de módulos do projeto}
\label{tab:hierarquia}
\begin{tabularx}{\linewidth}{llX}
\toprule
\textbf{Módulo} & \textbf{Arquivo} & \textbf{Função} \\
\midrule
\texttt{beat\_by\_bit\_de0\_cv} & wrapper board & Mapeia pinos físicos da DE0-CV \\
\texttt{proyecto}               & topo lógico   & Interliga todos os submódulos \\
\texttt{fluxo\_dados}           & FD            & Toda a lógica de jogo \\
\texttt{unidade\_controle}      & UC / FSM      & Máquina de estados do jogo \\
\texttt{packet\_sender}         & PS            & Serialização UART por quadro \\
\texttt{note\_track}            & NT (×4)       & Gerência de notas por trilha \\
\texttt{uart\_tx}               & TX            & Transmissor UART 8N1 \\
\texttt{debounce\_pulse}        & DB            & Debounce com saída em pulso \\
\texttt{contador\_m}            & CM            & Contador paramétrico com fim \\
\bottomrule
\end{tabularx}
\end{table}

% ==========================================================================
\section{Módulo \texttt{fluxo\_dados.v} — Expansão da ROM e Telemetria}
% ==========================================================================

O módulo \vfile{fluxo\_dados.v} é o núcleo do jogo. Ele concentra o gerador
de quadros, a ROM de chart, o contador de frames, os quatro \texttt{note\_track}
e os contadores de placar. Esta versão recebe duas alterações estruturais
significativas: a expansão da memória de chart e a exposição do frame counter
para telemetria.

% --------------------------------------------------------------------------
\subsection{Expansão da ROM de chart}

\subsubsection{Profundidade da ROM}

A constante \vsig{ROM\_DEPTH} foi aumentada de \antes{4428} para \depois{9861}
entradas. Cada entrada corresponde a um quadro de jogo a 60\,fps, portanto:

\begin{equation}
    t_{\max} = \frac{9861}{60} \approx 164{,}4\,\text{s}
    \qquad \text{(anteriormente } \frac{4428}{60} \approx 73{,}8\,\text{s)}
\end{equation}

Isso permite que o chart cubra a duração integral do arquivo WAV
\texttt{concerning\_hobbits.wav} sem truncamento.

\begin{lstlisting}[caption={Parâmetro ROM\_DEPTH — antes e depois}]
// Versao base
localparam integer ROM_DEPTH = 4428;

// Versao atual
localparam integer ROM_DEPTH = 9861;
\end{lstlisting}

\subsubsection{Largura das entradas}

Simultaneamente, a largura dos registros da ROM foi ampliada de 4 para 8 bits:

\begin{lstlisting}[caption={Tipo do array da ROM}]
// Versao base
reg [3:0] chart_rom [0:ROM_DEPTH-1];

// Versao atual
reg [7:0] chart_rom [0:ROM_DEPTH-1];
\end{lstlisting}

Os bits \texttt{[3:0]} continuam codificando o sinal de spawn para cada uma das
quatro trilhas. Os bits \texttt{[7:4]} estão reservados para extensões futuras
(ex.: notas com intensidade, notas longas).

% --------------------------------------------------------------------------
\subsection{Expansão do frame counter}

Com \vsig{ROM\_DEPTH} = 9\,861, o contador de frames precisa representar valores
até 9\,860, o que exige no mínimo $\lceil \log_2 9861 \rceil = 14$ bits
($2^{13} = 8192 < 9861 \leq 16383 = 2^{14}$). O registrador foi ampliado
de \antes{13} para \depois{14} bits:

\begin{lstlisting}[caption={Declaração do frame\_counter}]
// Versao base
reg [12:0] frame_counter;   // max 8191

// Versao atual
reg [13:0] frame_counter;   // max 16383 > 9861
\end{lstlisting}

A lógica de incremento e o teste de fim de chart (\vsig{fim\_tempo}) permanecem
funcionalmente idênticos; apenas os literais de inicialização foram atualizados
para o novo tamanho:

\begin{lstlisting}[caption={Reset e incremento do frame\_counter}]
// Versao base — reset assíncrono e síncrono agrupados no mesmo bloco
always @(posedge clock or posedge reset) begin
    if (reset || zera_timer) begin
        frame_counter <= 13'd0;
    end else if (frame_tick && game_active && frame_counter < ROM_DEPTH - 1) begin
        frame_counter <= frame_counter + 13'd1;
    end
end

// Versao atual — reset síncrono separado; tamanho 14 bits
always @(posedge clock or posedge reset) begin
    if (reset) begin
        frame_counter <= 14'd0;
    end else if (zera_timer) begin
        frame_counter <= 14'd0;
    end else if (frame_tick && game_active && frame_counter < ROM_DEPTH - 1) begin
        frame_counter <= frame_counter + 14'd1;
    end
end
\end{lstlisting}

\begin{destaque}
\textbf{Nota sobre o reset:} na versão base, \texttt{reset} e \texttt{zera\_timer}
compartilhavam o mesmo branch do \texttt{if}. Na versão atual eles foram separados
em branches distintos, o que é equivalente funcionalmente mas gera um código de
prioridade mais explícito e facilita a síntese correta.
\end{destaque}

% --------------------------------------------------------------------------
\subsection{Nova porta de saída: \texttt{frame\_counter\_out}}

Para permitir que o módulo de topo repasse a posição do chart ao
\texttt{packet\_sender}, foi adicionada uma porta de saída de 14 bits:

\begin{lstlisting}[caption={Nova porta de saída em fluxo\_dados}]
// Versao base — apenas frame_tick_out
output wire        frame_tick_out

// Versao atual — frame_tick_out + frame_counter_out
output wire        frame_tick_out,
output wire [13:0] frame_counter_out
\end{lstlisting}

A atribuição é trivial:

\begin{lstlisting}
assign frame_tick_out    = frame_tick;
assign frame_counter_out = frame_counter;   // linha adicionada
\end{lstlisting}

% --------------------------------------------------------------------------
\subsection{Correção no debounce do botão 0}

O debounce do botão 0 (\vsig{song\_ok\_pulso}) foi desacoplado do sinal
\vsig{limpaR} no pino de reset do submódulo \texttt{debounce\_pulse}:

\begin{lstlisting}[caption={Reset do debounce — botão 0}]
// Versao base
debounce_pulse #(.DEBOUNCE_CYCLES(BTN_DEBOUNCE_CYCLES))
    u_db_btn0 (.clock(clock), .reset(reset | limpaR), ...);

// Versao atual
debounce_pulse #(.DEBOUNCE_CYCLES(BTN_DEBOUNCE_CYCLES))
    u_db_btn0 (.clock(clock), .reset(reset), ...);
\end{lstlisting}

\subsubsection{Motivação}

A UC mantém \vsig{limpaR} em nível alto nos estados \texttt{idle} e
\texttt{select}. Na versão base, isso mantinha o debounce do botão 0
continuamente resetado enquanto o jogo estava em \texttt{select}, tornando
impossível detectar o pressionamento de seleção de música nesse estado.
Ao remover \vsig{limpaR} da condição de reset, o debounce opera normalmente
durante todo o ciclo — o que é seguro porque \vsig{song\_ok\_pulso} só é
interpretado pela UC nos estados \texttt{select} e \texttt{pause}.

% ==========================================================================
\section{Módulo \texttt{unidade\_controle.v} — Nova Transição no PAUSE}
% ==========================================================================

A Unidade de Controle (UC) implementa a FSM principal do jogo em Verilog
combinacional + registrador de estado. A versão atual introduz uma nova
transição de saída do estado \texttt{pause}.

% --------------------------------------------------------------------------
\subsection{Descrição dos estados}

\begin{table}[h!]
\centering
\caption{Estados da FSM}
\label{tab:estados}
\begin{tabular}{clp{8cm}}
\toprule
\textbf{Código} & \textbf{Estado} & \textbf{Descrição} \\
\midrule
\texttt{0000} & \texttt{idle}      & Tela inicial; contadores e placar limpos \\
\texttt{0001} & \texttt{countdown} & Contagem regressiva de 3 segundos \\
\texttt{0010} & \texttt{play}      & Jogo ativo; \texttt{game\_active} = 1 \\
\texttt{0011} & \texttt{pause}     & Jogo pausado \\
\texttt{0100} & \texttt{end\_win}  & Vitória (chart concluído sem 10 erros) \\
\texttt{0101} & \texttt{end\_lose} & Derrota (10 erros acumulados) \\
\texttt{0110} & \texttt{select}    & Seleção de música \\
\bottomrule
\end{tabular}
\end{table}

% --------------------------------------------------------------------------
\subsection{Nova transição: \texttt{PAUSE} $\to$ \texttt{IDLE}}

Na versão base, o único caminho de saída do estado \texttt{pause} era retornar
ao \texttt{play} via \vsig{start}. Na versão atual, foi adicionada uma
transição que retorna ao \texttt{idle} quando \vsig{song\_ok} (botão 0) é
pressionado:

\begin{lstlisting}[caption={Lógica de próximo estado — estado pause}]
// Versao base
pause: Eprox = start ? play : pause;

// Versao atual
pause: Eprox = song_ok ? idle :
               start   ? play  : pause;
\end{lstlisting}

\subsubsection{Diagrama de estados atualizado}

\begin{figure}[h!]
\centering
\begin{tikzpicture}[
    ->, >=Stealth,
    node distance=2.2cm and 3.2cm,
    state/.style={circle, draw, minimum size=1.5cm, font=\small\bfseries,
                  fill=white, draw=sectioncolor, line width=1pt, align=center},
    every label/.style={font=\scriptsize}
]

% Nós
\node[state] (idle)      {idle};
\node[state, right=of idle] (select)    {select};
\node[state, right=of select] (countdown) {count\\down};
\node[state, right=of countdown] (play) {play};
\node[state, below=of play]  (pause)    {pause};
\node[state, below=of countdown] (lose) {end\\lose};
\node[state, below=of select] (win)     {end\\win};

% Transições normais
\draw (idle)      edge[bend left=10]  node[above,font=\scriptsize]{\textit{start}} (select);
\draw (select)    edge[bend left=10]  node[above,font=\scriptsize]{\textit{song\_ok}} (countdown);
\draw (select)    edge[bend left=15]  node[below,font=\scriptsize]{\textit{start}} (idle);
\draw (countdown) edge[bend left=10]  node[above,font=\scriptsize]{\textit{fim\_contagem}} (play);
\draw (play)      edge[bend left=10]  node[right,font=\scriptsize]{\textit{start}} (pause);
\draw (pause)     edge[bend left=10]  node[left,font=\scriptsize]{\textit{start}} (play);
\draw (play)      edge[bend left=5]   node[left,font=\scriptsize]{\textit{perdeu}} (lose);
\draw (play)      edge[bend left=15]  node[above,font=\scriptsize]{\textit{fim\_tempo}} (win);
\draw (win)       edge[bend left=10]  node[left,font=\scriptsize]{\textit{start}} (idle);
\draw (lose)      edge[bend left=15]  node[below,font=\scriptsize]{\textit{start}} (idle);

% Nova transição destacada
\draw[color=newcolor, line width=1.5pt, dashed]
    (pause) to[bend right=35]
    node[below,font=\scriptsize\bfseries,color=newcolor]{\textit{song\_ok} (NOVO)}
    (idle);

% Self-loops
\draw (idle)      edge[loop left]  node[font=\scriptsize]{} (idle);
\draw (select)    edge[loop above] node[font=\scriptsize]{} (select);
\draw (pause)     edge[loop right] node[font=\scriptsize]{} (pause);
\draw (countdown) edge[loop above] node[font=\scriptsize]{} (countdown);
\draw (win)       edge[loop left]  node[font=\scriptsize]{} (win);
\draw (lose)      edge[loop below] node[font=\scriptsize]{} (lose);

\end{tikzpicture}
\caption{FSM do jogo — transição nova destacada em verde tracejado}
\label{fig:fsm}
\end{figure}

\subsubsection{Impacto na experiência do usuário}

Antes da mudança, abandonar uma música pausada exigia pressionar o botão de
reset físico da placa. Agora, pressionar o botão 0 (\texttt{song\_ok}) em
\texttt{pause} retorna ao \texttt{idle}, de onde o fluxo segue normalmente:
\texttt{idle} $\to$ \texttt{select} $\to$ \texttt{countdown} $\to$
\texttt{play}. O placar é zerado automaticamente porque os estados
\texttt{idle} e \texttt{select} mantêm \vsig{limpaR} = 1.

% --------------------------------------------------------------------------
\subsection{Saídas da UC por estado}

\begin{table}[h!]
\centering
\caption{Tabela de saídas da UC por estado}
\label{tab:saidas-uc}
\small
\begin{tabular}{lcccccc}
\toprule
\textbf{Estado}    & \vsig{limpaR} & \vsig{zera\_timer} & \vsig{conta\_timer} & \vsig{game\_active} & \vsig{registraR} \\
\midrule
\texttt{idle}      & 1 & 1 & 0 & 0 & 0 \\
\texttt{select}    & 1 & 1 & 0 & 0 & 0 \\
\texttt{countdown} & 0 & 0 & 1 & 0 & 0 \\
\texttt{play}      & 0 & 0 & 1 & 1 & \textit{jogada} \\
\texttt{pause}     & 0 & 0 & 0 & 0 & 0 \\
\texttt{end\_win}  & 0 & 1 & 0 & 0 & 0 \\
\texttt{end\_lose} & 0 & 1 & 0 & 0 & 0 \\
\bottomrule
\end{tabular}
\end{table}

% ==========================================================================
\section{Módulo \texttt{packet\_sender.v} — Telemetria de Posição via UART}
% ==========================================================================

O \vfile{packet\_sender.v} monta e envia um pacote de bytes pela UART a cada
\vsig{frame\_tick} ($\approx$60 pacotes/segundo). A versão atual adiciona o
frame counter ao pacote, expandindo-o de 20 para 22 bytes.

% --------------------------------------------------------------------------
\subsection{Nova porta de entrada}

\begin{lstlisting}[caption={Declaração de porta adicionada ao packet\_sender}]
// Versao base — sem frame_counter
// Versao atual
input wire [13:0] frame_counter,    // progresso do chart
\end{lstlisting}

% --------------------------------------------------------------------------
\subsection{Formato do pacote UART}

A Tabela~\ref{tab:pacote} detalha o formato completo dos dois pacotes.

\begin{table}[h!]
\centering
\caption{Formato do pacote UART — comparação entre versões}
\label{tab:pacote}
\small
\begin{tabular}{cllcc}
\toprule
\textbf{Byte} & \textbf{Campo} & \textbf{Valor / Sinal} & \textbf{Base} & \textbf{Atual} \\
\midrule
\texttt{[0]}  & Marcador início & \texttt{0x55}                     & \checkmark & \checkmark \\
\texttt{[1]}  & Estado do jogo  & \texttt{\{4'b0, game\_state\}}    & \checkmark & \checkmark \\
\texttt{[2]}  & Countdown       & \texttt{countdown\_sec}            & \checkmark & \checkmark \\
\texttt{[3]}  & T0 nota 0       & \texttt{t0n0} (age; 0xFF=vazio)   & \checkmark & \checkmark \\
\texttt{[4]}  & T0 nota 1       & \texttt{t0n1}                     & \checkmark & \checkmark \\
\texttt{[5]}  & T0 nota 2       & \texttt{t0n2}                     & \checkmark & \checkmark \\
\texttt{[6]}  & T1 nota 0       & \texttt{t1n0}                     & \checkmark & \checkmark \\
\texttt{[7]}  & T1 nota 1       & \texttt{t1n1}                     & \checkmark & \checkmark \\
\texttt{[8]}  & T1 nota 2       & \texttt{t1n2}                     & \checkmark & \checkmark \\
\texttt{[9]}  & T2 nota 0       & \texttt{t2n0}                     & \checkmark & \checkmark \\
\texttt{[10]} & T2 nota 1       & \texttt{t2n1}                     & \checkmark & \checkmark \\
\texttt{[11]} & T2 nota 2       & \texttt{t2n2}                     & \checkmark & \checkmark \\
\texttt{[12]} & T3 nota 0       & \texttt{t3n0}                     & \checkmark & \checkmark \\
\texttt{[13]} & T3 nota 1       & \texttt{t3n1}                     & \checkmark & \checkmark \\
\texttt{[14]} & T3 nota 2       & \texttt{t3n2}                     & \checkmark & \checkmark \\
\texttt{[15]} & Score alto      & \texttt{score[15:8]}              & \checkmark & \checkmark \\
\texttt{[16]} & Score baixo     & \texttt{score[7:0]}               & \checkmark & \checkmark \\
\texttt{[17]} & Misses          & \texttt{misses}                   & \checkmark & \checkmark \\
\texttt{[18]} & Combo           & \texttt{combo}                    & \checkmark & \checkmark \\
\texttt{[19]} & Frame alto      & \texttt{\{2'b00, frame\_counter[13:8]\}} & \texttt{0xAA}$^*$ & \checkmark \\
\texttt{[20]} & Frame baixo     & \texttt{frame\_counter[7:0]}      & ---        & \checkmark \\
\texttt{[21]} & Marcador fim    & \texttt{0xAA}                     & ---        & \checkmark \\
\midrule
\multicolumn{3}{l}{\textbf{Total}} & \textbf{20\,B} & \textbf{22\,B} \\
\bottomrule
\multicolumn{5}{l}{$^*$ Na versão base, o marcador \texttt{0xAA} ocupava o byte \texttt{[19]}.}
\end{tabular}
\end{table}

% --------------------------------------------------------------------------
\subsection{Alterações no buffer e na condição de término}

\begin{lstlisting}[caption={Buffer do pacote}]
// Versao base
reg [7:0] pkt [0:19];   // 20 bytes

// Versao atual
reg [7:0] pkt [0:21];   // 22 bytes
\end{lstlisting}

\begin{lstlisting}[caption={Preenchimento do pacote no estado WAIT}]
// Versao base
pkt[19] <= 8'hAA;

// Versao atual
pkt[19] <= {2'b00, frame_counter[13:8]};
pkt[20] <= frame_counter[7:0];
pkt[21] <= 8'hAA;
\end{lstlisting}

\begin{lstlisting}[caption={Condição de último byte no estado SEND}]
// Versao base
if (byte_idx == 19) begin
    state <= DONE;
end

// Versao atual
if (byte_idx == 21) begin
    state <= DONE;
end
\end{lstlisting}

% --------------------------------------------------------------------------
\subsection{Latência do pacote}

Com 22 bytes a 115\,200 bps (10 bits por byte em modo 8N1):

\begin{equation}
    t_{\text{pacote}} = \frac{22 \times 10}{115200} \approx 1{,}91\,\text{ms}
    \qquad \Bigl(\text{período de quadro} = \frac{1}{60} \approx 16{,}67\,\text{ms}\Bigr)
\end{equation}

O tempo de transmissão representa apenas 11,4\% do período de quadro,
não havendo risco de o \texttt{packet\_sender} não terminar antes do
próximo \vsig{frame\_tick}.

% --------------------------------------------------------------------------
\subsection{Decodificação do frame counter no host}

O Python pode reconstruir a posição do chart a partir dos dois bytes
transmitidos com a seguinte operação:

\begin{equation}
    \text{frame} = (\texttt{pkt[19]} \,\&\, \texttt{0x3F}) \ll 8 \;\mid\; \texttt{pkt[20]}
\end{equation}

Isso permite sincronizar o início da reprodução de áudio com o frame exato
em que o chart se encontra, eliminando desalinhamentos acumulados por
variações de clock entre o PC e a FPGA.

% ==========================================================================
\section{Módulo de Topo \texttt{proyecto.v} — Interligação}
% ==========================================================================

O módulo de topo conecta \texttt{fluxo\_dados}, \texttt{unidade\_controle} e
\texttt{packet\_sender}. As alterações aqui são puramente de cabeamento,
decorrentes das novas portas adicionadas nos módulos internos.

% --------------------------------------------------------------------------
\subsection{Novo fio de interligação}

\begin{lstlisting}[caption={Declaração do fio s\_frame\_counter}]
// Versao base — nao existe
// Versao atual
wire [13:0] s_frame_counter;
\end{lstlisting}

% --------------------------------------------------------------------------
\subsection{Atualização nas instâncias}

\begin{lstlisting}[caption={Instância de fluxo\_dados — nova conexão}]
fluxo_dados u_fd (
    ...
    .frame_tick_out    (s_frame_tick),
    .frame_counter_out (s_frame_counter)   // NOVO
);
\end{lstlisting}

\begin{lstlisting}[caption={Instância de packet\_sender — nova conexão}]
packet_sender u_ps (
    ...
    .combo         (s_combo),
    .frame_counter (s_frame_counter),      // NOVO
    .uart_tx_out   (uart_tx_out)
);
\end{lstlisting}

% ==========================================================================
\section{Arquivo de Constraints \texttt{beat\_by\_bit.sdc}}
% ==========================================================================

O único ponto alterado no arquivo de timing constraints é a referência ao
pino de clock:

\begin{lstlisting}[language={},caption={Constraint de clock}]
# Versao base
create_clock -name clk -period 20.000 [get_ports {CLOCK_50}]

# Versao atual
create_clock -name clk -period 20.000 [get_ports {clock}]
\end{lstlisting}

A mudança reflete a renomeação interna do pino no módulo wrapper
\texttt{beat\_by\_bit\_de0\_cv.v}, que mapeia o sinal \texttt{CLOCK\_50}
da placa para o nome genérico \texttt{clock}. O período de 20\,ns
corresponde ao clock de 50\,MHz e permanece inalterado.

% ==========================================================================
\section{Módulos Sem Alteração}
% ==========================================================================

Os módulos listados abaixo permaneceram funcionalmente idênticos entre as
duas versões. Nenhuma porta, parâmetro ou lógica interna foi modificada.

\begin{table}[h!]
\centering
\caption{Módulos inalterados}
\begin{tabular}{ll}
\toprule
\textbf{Arquivo} & \textbf{Função} \\
\midrule
\texttt{uart\_tx.v}          & Transmissor UART 8N1 paramétrico \\
\texttt{note\_track.v}       & Gerência de notas por trilha (até 3 simultâneas) \\
\texttt{edge\_detector.v}    & Detector de borda de subida/descida \\
\texttt{debounce\_pulse.v}   & Debounce com saída em pulso \\
\texttt{hexa7seg.v}          & Decodificador para display 7 segmentos \\
\texttt{lfsr8.v}             & LFSR de 8 bits \\
\texttt{registrador\_4.v}    & Registrador de 4 bits com enable e clear \\
\texttt{contador\_m.v}       & Contador paramétrico M com saída de fim \\
\texttt{beat\_by\_bit\_de0\_cv.v} & Wrapper de pinos da placa DE0-CV \\
\bottomrule
\end{tabular}
\end{table}

% ==========================================================================
\section{Tabela Consolidada de Mudanças}
% ==========================================================================

\begin{longtable}{p{3.5cm} p{5cm} p{5cm}}
\caption{Resumo consolidado de todas as alterações} \label{tab:consolidado} \\
\toprule
\textbf{Arquivo} & \textbf{Versão base} & \textbf{Versão atual} \\
\midrule
\endfirsthead
\toprule
\textbf{Arquivo} & \textbf{Versão base} & \textbf{Versão atual} \\
\midrule
\endhead
\bottomrule
\endfoot

\texttt{fluxo\_dados.v} &
    \texttt{ROM\_DEPTH = 4428} &
    \texttt{ROM\_DEPTH = 9861} \\[4pt]

\texttt{fluxo\_dados.v} &
    \texttt{reg [3:0] chart\_rom} &
    \texttt{reg [7:0] chart\_rom} \\[4pt]

\texttt{fluxo\_dados.v} &
    \texttt{reg [12:0] frame\_counter} &
    \texttt{reg [13:0] frame\_counter} \\[4pt]

\texttt{fluxo\_dados.v} &
    Sem porta \texttt{frame\_counter\_out} &
    \texttt{output wire [13:0] frame\_counter\_out} \\[4pt]

\texttt{fluxo\_dados.v} &
    \texttt{debounce btn0:} \newline \texttt{.reset(reset | limpaR)} &
    \texttt{debounce btn0:} \newline \texttt{.reset(reset)} \\[4pt]

\texttt{unidade\_controle.v} &
    \texttt{pause}: \newline $\text{start} \to \text{play}$ &
    \texttt{pause}: \newline $\text{song\_ok} \to \text{idle}$, $\text{start} \to \text{play}$ \\[4pt]

\texttt{packet\_sender.v} &
    Sem porta \texttt{frame\_counter} &
    \texttt{input wire [13:0] frame\_counter} \\[4pt]

\texttt{packet\_sender.v} &
    \texttt{pkt[0:19]} (20\,B) &
    \texttt{pkt[0:21]} (22\,B) \\[4pt]

\texttt{packet\_sender.v} &
    \texttt{pkt[19] = 0xAA} &
    \texttt{pkt[19..20] = frame\_counter} \newline \texttt{pkt[21] = 0xAA} \\[4pt]

\texttt{packet\_sender.v} &
    \texttt{byte\_idx == 19} &
    \texttt{byte\_idx == 21} \\[4pt]

\texttt{proyecto.v} &
    Sem \texttt{s\_frame\_counter} &
    \texttt{wire [13:0] s\_frame\_counter} \\[4pt]

\texttt{beat\_by\_bit.sdc} &
    \texttt{get\_ports \{CLOCK\_50\}} &
    \texttt{get\_ports \{clock\}} \\

\end{longtable}

% ==========================================================================
\section{Impacto nas Capacidades do Sistema}
% ==========================================================================

\begin{table}[h!]
\centering
\caption{Comparativo de capacidades}
\label{tab:capacidades}
\begin{tabular}{lcc}
\toprule
\textbf{Capacidade} & \textbf{Versão base} & \textbf{Versão atual} \\
\midrule
Duração máxima do chart      & $\approx$73\,s    & $\approx$164\,s \\
Bits do frame counter        & 13\,bits          & 14\,bits \\
Largura da ROM               & 4\,bits/frame     & 8\,bits/frame \\
Tamanho do pacote UART       & 20\,bytes         & 22\,bytes \\
Tempo de tx por pacote       & $\approx$1,74\,ms & $\approx$1,91\,ms \\
Telemetria de posição        & Não               & Sim \\
Saída do PAUSE sem reset     & Não               & Sim (via botão 0) \\
\bottomrule
\end{tabular}
\end{table}

\vspace{1em}

\begin{destaque}
\textbf{Compatibilidade de interface:} a interface de I/O externa do sistema
(pinos físicos da DE0-CV, protocolo UART de nível físico) permanece idêntica.
As mudanças são internas ao projeto RTL e ao protocolo de dados do pacote.
O software Python precisa ser atualizado para ler 22 bytes por pacote (em vez
de 20) e para decodificar os bytes \texttt{[19]} e \texttt{[20]} como
\texttt{frame\_counter}.
\end{destaque}

\end{document}
